% ======================================================================
%  Solution slides TEMPLATE: an empty deck in the shared style.
%  Copied to slides/slides.tex; fill it in.
%  - WHICH slides come in WHICH order and WHAT they say is defined in
%    pattern.md (given to you in the prompt).  This file only shows HOW
%    each kind of slide is laid out: one frame per kind.
%  - Repeat the frames per concept / problem / task, delete the ones you
%    do not need, delete the layout comments when a frame is filled.
%  - Replace every X.Y (problem) and X.Y.Z (task) and every title.
%  - Keep the stage label (\Refresh, \Given, \Byhand, \Algo, \Impl, \Answer) as the
%    first line of each frame.
%  - Do not change the preamble or define new colours/macros: the look
%    lives in exslides.sty and must be identical for every sheet.
% ======================================================================
\documentclass[aspectratio=169,11pt]{beamer}
\usepackage{exslides}

\title{Exercise X -- Problems X.Y and X.Y}
\subtitle{Course (course number), term}
\date{Solution slides}
\begin{document}
\maketitle

% ======================================================================
\section{Refresher: what you should remember}

\begin{frame}{Refresher: concept}
\Refresh
% one sentence, then \begin{enumerate} with display formulas \[...\];
% a figure may sit centred above a display formula
\end{frame}

% ======================================================================
\section{Problem X.Y: recap of the task}

\begin{frame}{Given}
\Given
% \begin{itemize} ... \end{itemize}
% with an image from the sheet: \begin{columns}[T] \column{.42\textwidth} list
% \column{.56\textwidth} \includegraphics[width=\linewidth]{sheet_pP_NAME.png}
% {\footnotesize sheet, p.~P} \end{columns}
\end{frame}

\begin{frame}{Overview: the tasks}\small
\Given
% \begin{description}[Task X.Y.Z] \item[Task X.Y.Z] short title: a few words \end{description}
\end{frame}

% ======================================================================
\section{Task X.Y.Z: short task title}

\begin{frame}{Background: question it answers}
\Given
% \begin{itemize} short steps \end{itemize}, display formulas \[ ... \quad(5) \]
% with the sheet's equation numbers; last item: \item[$\Rightarrow$] what it means;
% with an image: two columns as on "Given"
\end{frame}

\begin{frame}{X.Y.Z -- what is asked}\small
\Given
% the task restated; display formulas with equation numbers; then \medskip
% and the sheet's warning or expectation; with an image: two columns as on "Given"
\end{frame}

\begin{frame}{X.Y.Z -- by hand on a toy example}\small
\Byhand
% the toy data as a display formula, \begin{enumerate} of steps,
% then \medskip and \alert{Why ...?}
% (only if the real data is too big for hand work or hides the point)
\end{frame}

\begin{frame}{X.Y.Z -- by hand: the real data}
\Byhand
% \begin{center}\begin{tabular}{lr}\toprule ... \bottomrule\end{tabular}\end{center}
% then \medskip and "Sanity checks: ..."
\end{frame}

\begin{frame}{X.Y.Z -- algorithm}
\Algo
% \begin{enumerate} of steps, then \medskip and \alert{Pitfall:}
\end{frame}

\begin{frame}[fragile]{X.Y.Z -- implementation}
\Impl
% \begin{lstlisting} ... \end{lstlisting} (frame must be [fragile]);
% with a figure: \begin{columns}[T] \column{.55\textwidth} code
% \column{.43\textwidth} \includegraphics[width=\linewidth]{file.png}
% \small \begin{itemize} notes \end{itemize} \end{columns}
\end{frame}

\begin{frame}{X.Y.Z -- result and checks}
\Impl
% \[\boxed{...}\], then \textbf{Checks} and \begin{itemize}
\end{frame}

\begin{frame}{X.Y.Z -- what the plot shows}
\Impl
% \includegraphics[width=\linewidth]{file.png}, then \small one or two sentences
\end{frame}

% --- task of type "question list": one frame per question instead of the steps
\begin{frame}{X.Y.Z -- question 1: short title}
\Answer
% \textbf{Question 1} (or the sheet's label) the question as on the sheet; \medskip then the answer in one
% sentence, then the reason / formula / example (at most three sentences)
\end{frame}

\begin{frame}{Why is the result not the same as in X.Y?}
% no stage label; \begin{center} booktabs table \end{center}, then one sentence
\end{frame}

% ======================================================================
\section{Takeaways}

\begin{frame}{Takeaways}\small
% \begin{itemize} one line per lesson
\end{frame}

\end{document}
